\chapter{البنية الداخلية وطريقة عمل محركات السيرفو}
\label{ch:servo_anatomy}

% مقدمة تمهيدية للفصل
بعد أن استعرضنا الهيكلية العامة لأنظمة التحكم بالحركة في الفصل الأول، ننتقل في هذا الفصل إلى دراسة القلب النابض لهذه الأنظمة، وهو محرك السيرفو (\textenglish{Servo Motor}). يتميز هذا المحرك عن المحركات الكهربائية التقليدية بقدرته الفائقة على الاستجابة اللحظية للأوامر، وتوفير عزم دوران عالٍ عند سرعات مختلفة. سيتناول هذا الفصل البنية الفيزيائية الداخلية للمحرك، وأنظمة التغذية الراجعة المدمجة به، وصولاً إلى مقارنة عملية بين محركات السيرفو التي تعمل بالتيار المستمر وتلك التي تعمل بالتيار المتناوب.

\vspace{0.5cm}

% استدعاء ملفات السكشنات
\section{التركيب الفيزيائي لمحركات السيرفو}
\label{sec:physical_structure}

تعتمد محركات السيرفو الحديثة، وتحديداً المحركات التزامنية ذات المغناطيس الدائم (\textenglish{PMSM - Permanent Magnet Synchronous Motors})، على بنية هندسية دقيقة تجمع بين الكفاءة الكهرومغناطيسية وانخفاض عزم القصور الذاتي الميكانيكي (\textenglish{Low Inertia}). هذا التصميم المدمج هو ما يمنحها القدرة على التسارع والتباطؤ بشكل شبه لحظي.

يتكون محرك السيرفو من جزأين فيزيائيين رئيسيين:

\subsection{العضو الثابت (\textenglish{Stator})}
يمثل الهيكل الخارجي الثابت للمحرك، ويتألف من مجموعة من الرقائق المعدنية (\textenglish{Laminations}) المصنوعة من صلب السليكون، وذلك لتقليل المفاقيد المغناطيسية الناتجة عن التيارات الدوامية (\textenglish{Eddy Currents}). تُلف أسلاك نحاسية معزولة داخل مجاري العضو الثابت لتشكيل ملفات ثلاثية الطور (\textenglish{3-Phase Windings}). عند مرور تيار كهربائي قادم من المشغل (\textenglish{Servo Drive}) في هذه الملفات، يتولد مجال مغناطيسي دوار يعبر الفجوة الهوائية باتجاه العضو الدوار.

\subsection{العضو الدوار (\textenglish{Rotor})}
على عكس المحركات الحثية التقليدية التي تعتمد على تحريض التيار في العضو الدوار، يُصنع العضو الدوار في محركات السيرفو الحديثة من مغانط دائمة قوية جداً من مواد أرضية نادرة، مثل النيوديميوم (\textenglish{Neodymium}). يحقق هذا التصميم ميزتين جوهريتين:
\begin{itemize}
    \item \textbf{غياب المفاقيد النحاسية في العضو الدوار:} لعدم وجود ملفات نحاسية فيه، مما يقلل بشكل كبير من ارتفاع حرارة المحرك ويزيد من كفاءته.
    \item \textbf{انخفاض عزم القصور الذاتي:} نظراً لصغر حجم وخفة وزن المغانط الدائمة مقارنة بالملفات التقليدية، مما يسمح للمحرك بالاستجابة الديناميكية السريعة جداً لتغيرات السرعة والاتجاه المطلوبة من المتحكم.
\end{itemize}

\subsection{مبدأ توليد العزم الكهرومغناطيسي}
ينشأ العزم الميكانيكي (\textenglish{Torque}) نتيجة التجاذب والتنافر المستمر بين المجال المغناطيسي الدوار الناشئ في العضو الثابت، والمجال المغناطيسي الثابت للمغانط في العضو الدوار. يمكن التعبير عن العزم الميكانيكي المتولد رياضياً باستخدام المعادلة المبسطة التالية:
$$T = K_t \cdot I_q$$
حيث:
\begin{itemize}
    \item $T$: العزم الميكانيكي المتولد بوحدة نيوتن-متر (\textenglish{N.m}).
    \item $K_t$: ثابت العزم الخاص بالمحرك (\textenglish{Torque Constant})، وهو قيمة فيزيائية تحددها الشركة المصنعة بناءً على قوة المغانط وعدد لفات النحاس.
    \item $I_q$: التيار الفعال (العمودي) المار في ملفات العضو الثابت، والمتحكم فعلياً في توليد العزم.
\end{itemize}
\section{أنظمة التغذية الراجعة (\textenglish{Feedback Systems})}
\label{sec:feedback_systems}

السمة الجوهرية التي تُميز محرك السيرفو عن بقية المحركات الكهربائية التقليدية هي وجود حساس تغذية راجعة مُدمج، يُثبت عادةً في الجزء الخلفي من عمود الدوران (\textenglish{Motor Shaft}). تكمن وظيفة هذا الحساس في قراءة الموضع الزاوي، السرعة، واتجاه الدوران بشكل لحظي، ثم إرسال هذه البيانات إلى وحدة التحكم لإغلاق حلقة التحكم وتصحيح مسار الحركة.

تُقسم أنظمة التغذية الراجعة المستخدمة في الصناعة إلى عدة أنواع بناءً على تقنية القياس ومستوى الدقة المطلوب:

\subsection{المشفرات (\textenglish{Encoders})}
تُعد المشفرات (الضوئية أو المغناطيسية) الحساسات الأكثر شيوعاً في تطبيقات الأتمتة الدقيقة. وتعتمد على تحويل الحركة الميكانيكية إلى إشارات رقمية يفهمها المتحكم. تنقسم المشفرات إلى فئتين أساسيتين:

\begin{enumerate}
    \item \textbf{المشفرات التزايدية أو النسبية (\textenglish{Incremental Encoders}):}
    يقوم هذا النوع بتوليد سلسلة من النبضات الكهربائية (\textenglish{Pulses}) المتساوية مع كل زيادة في زاوية الدوران. يقوم المتحكم بحساب الموضع الفعلي عن طريق عَدّ هذه النبضات انطلاقاً من نقطة مرجعية. 
    \begin{itemize}
        \item \textbf{الميزة:} بساطة التصميم وانخفاض التكلفة.
        \item \textbf{العيب:} عند انقطاع التيار الكهربائي، يفقد النظام قراءة الموضع الحالي، مما يستدعي إجراء عملية "العودة إلى نقطة الصفر" (\textenglish{Homing Procedure}) عند إعادة التشغيل لتعريف الآلة بمكانها الصحيح.
    \end{itemize}
    
    \item \textbf{المشفرات المطلقة (\textenglish{Absolute Encoders}):}
    على عكس النوع التزايدي، يقرأ هذا المشفر الموقع باستخدام قرص مشفر يحمل نمطاً معقداً، بحيث يقابل كل موضع زاوي كوداً رقمياً فريداً (مثل \textenglish{Gray Code}).
    \begin{itemize}
        \item \textbf{الميزة الجوهرية:} يحتفظ بقيمة الموضع الفعلي والدقيق حتى في حال انقطاع الطاقة الكهربائية، مما يجعله الخيار القياسي والآمن في التطبيقات الحرجة كالأذرع الروبوتية، حيث قد تتسبب عملية الـ \textenglish{Homing} باصطدام الروبوت بالمعدات المحيطة.
    \end{itemize}
\end{enumerate}

\subsection{المحللات أو الريزولفرات (\textenglish{Resolvers})}
المحلل هو عبارة عن محول كهربائي دوار (\textenglish{Rotary Electrical Transformer}) يُستخدم لقياس زاوية الدوران. على عكس المشفرات التي تعطي إشارات رقمية، يقوم المحلل بتوليد إشارات تماثلية (\textenglish{Analog Signals}) متمثلة في موجات جيبية وجيب تمام (\textenglish{Sine and Cosine}) تتناسب سعتها مع زاوية العضو الدوار.

رغم أن المحللات قد توفر دقة أقل نسبياً من المشفرات الضوئية المتقدمة، إلا أنها تتفوق في جانب بالغ الأهمية: \textbf{المتانة الميكانيكية والاعتمادية}. لكونها تعتمد على ملفات نحاسية وقلب حديدي (بدون دوائر إلكترونية دقيقة أو أقراص زجاجية قابلة للكسر)، فإنها تتحمل ظروف العمل الصناعية القاسية جداً، مثل درجات الحرارة المرتفعة، والاهتزازات العنيفة، والأتربة، مما يجعلها مفضلة في الصناعات الثقيلة والتطبيقات العسكرية.
\section{مقارنة بين محركات السيرفو للتيار المستمر والمتناوب}
\label{sec:ac_dc_comparison}

تاريخياً، كانت محركات التيار المستمر (\textenglish{DC Servo Motors}) هي الخيار السائد في أنظمة التحكم نظراً لسهولة نمذجتها الرياضية وبساطة دوائر التحكم في سرعتها وعزمها. ولكن مع التطور الهائل في إلكترونيات القدرة (\textenglish{Power Electronics})، وتقنيات المعالجات الدقيقة، أصبحت محركات التيار المتناوب (\textenglish{AC Servo Motors}) هي المعيار الذهبي المعتمد في الأتمتة الصناعية الحديثة.

\subsection{خصائص سيرفو التيار المستمر (\textenglish{DC Servo})}
تُستخدم هذه المحركات غالباً في التطبيقات ذات الاستطاعات المنخفضة إلى المتوسطة، أو في الأنظمة التي تعتمد على البطاريات (مثل الروبوتات المتنقلة). وتعتمد إما على المحركات ذات الفرش (\textenglish{Brushed DC}) أو بدون فرش (\textenglish{BLDC}).
\begin{itemize}
    \item \textbf{المزايا:} سهولة التحكم والمناورة، وتكلفة ابتدائية أقل لمشغلات القدرة (\textenglish{Drives}).
    \item \textbf{القيود:} في الأنواع ذات الفرش الكربونية، يتطلب الأمر صيانة دورية بسبب التآكل الميكانيكي للفرش. كما أن كفاءة تبديد الحرارة تكون أقل نظراً لتولد الحرارة في العضو الدوار الداخلي.
\end{itemize}

\subsection{خصائص سيرفو التيار المتناوب (\textenglish{AC Servo})}
تعتمد بشكل أساسي على المحركات التزامنية ذات المغناطيس الدائم (\textenglish{PMSM}). وهي الخيار الأول والتطبيق الفعلي للروبوتات الصناعية وآلات القطع المبرمجة (\textenglish{CNC}).
\begin{itemize}
    \item \textbf{المزايا:} غياب التآكل الميكانيكي للفرش (انعدام الحاجة للصيانة الدورية)، وكفاءة تبريد عالية جداً؛ حيث تقع الملفات النحاسية على العضو الثابت الخارجي، مما يسهل تشتيت الحرارة إلى الهيكل المعدني للآلة. بالإضافة إلى توفير عزم مستقر حتى عند السرعات العالية جداً.
    \item \textbf{القيود:} تتطلب خوارزميات تحكم رياضية معقدة جداً (مثل التحكم الموجه للمجال \textenglish{Field Oriented Control - FOC})، وتكلفتها الإجمالية أعلى مقارنة بنظيرتها في التيار المستمر.
\end{itemize}

\subsection{مقارنة شاملة}
يلخص الجدول (\ref{tab:ac_dc_comparison}) أبرز الفروقات التشغيلية والهندسية بين النوعين:

\begin{table}[h!]
    \centering
    \renewcommand{\arraystretch}{1.4}
    \begin{tabular}{|p{3.5cm}|p{5cm}|p{5cm}|}
        \hline
        \textbf{وجه المقارنة} & \textbf{سيرفو التيار المستمر (\textenglish{DC})} & \textbf{سيرفو التيار المتناوب (\textenglish{AC})} \\
        \hline
        \textbf{الصيانة والعمر الافتراضي} & تتطلب صيانة (لتبديل الفرش) وعمرها الافتراضي محدود & خالية من الصيانة وعمرها الافتراضي طويل جداً \\
        \hline
        \textbf{تبديد الحرارة} & ضعيف نسبياً & ممتاز (يسهل تبريد العضو الثابت الخارجي) \\
        \hline
        \textbf{أداء السرعة والعزم} & ينخفض العزم المتاح تدريجياً مع ارتفاع السرعة & عزم ثابت ومستقر عبر نطاق واسع جداً من السرعات \\
        \hline
        \textbf{التطبيقات الشائعة} & الروبوتات المتنقلة (\textenglish{AGVs})، والمعدات الطبية، والأجهزة الصغيرة & الأذرع الروبوتية ذات الدرجات المتعددة، آلات \textenglish{CNC}، وخطوط التعبئة السريعة \\
        \hline
    \end{tabular}
    \caption[مقارنة هندسية وعملية بين محركات السيرفو للتيار المستمر والتيار المتناوب]{\centering مقارنة هندسية وعملية بين محركات السيرفو للتيار المستمر والتيار المتناوب}
    \label{tab:ac_dc_comparison}
\end{table}